\section{Benchmark Coverage and Validity}
\label{sec:benchmarks}
A graphic-design benchmark is only informative relative to its \emph{input contract}, \emph{output contract}, \emph{task}, and \emph{scoring target}.
Table~\ref{tab:benchmarks} separates generation, visual judgment, and editing.
Their numerical results must not be aggregated into a cross-task leaderboard without calibration.

% Independent source table.
\begin{table*}[t]
\centering
\caption{Representative design benchmarks and limits of their evaluation contracts.
The table identifies studied tasks; it does not claim that every metric has been independently reproduced.}
\label{tab:benchmarks}
\footnotesize
\begin{tabularx}{\textwidth}{@{}p{0.19\textwidth}p{0.17\textwidth}p{0.25\textwidth}X@{}}
\toprule
\textbf{Benchmark} & \textbf{Artifact / task} & \textbf{Assessed capability} & \textbf{Evidence boundary}\\
\midrule
SVGEditBench \citep{svgeditbench2024} & SVG code editing & Prompt-directed SVG modifications & Does not establish native edit quality in raster/PSD/DOM.\\
VectorGym \citep{vectorgym2026} & SVG; text/sketch/code & SVG generation, editing, sketch conversion, captioning & Four tasks and their metrics are not interchangeable.\\
MiLDEdit \citep{mildedit2026} & Multi-layer design documents & Instruction, layout, aesthetics, text & Abstract-level evidence for the four axes does not establish universal scoring.\\
InfoEdit \citep{yang2026infoedit} & Infographics; edit requests & Semantic edit and layout \emph{reflow} & Rigidly unchanged non-target pixels can be an incorrect target.\\
MPTc-Bench \citep{lin2026mptc} & Movie-poster adaptation & Film-identity retention; typography; target-market fit & Market adaptation changes content in ways pure similarity misses.\\
GraphicDesignBench \citep{deganutti2026gdb} & Layered, SVG, text, animation & Spatial accuracy, text fidelity, semantic/structural validity & Task counts depend on the version of paper and project.\\
BizGenEval \citep{li2026bizgeneval} & Rendered commercial visuals & Layout, typography, attributes, knowledge checklists & Successful raster rendering does not guarantee editable source objects.\\
AesEval-Bench \citep{an2026aeseval} & Graphic-design aesthetic assessment & Aesthetic judgment, region selection, localization & Judge performance is not design editing performance.\\
\bottomrule
\end{tabularx}
\end{table*}


\subsection{What a Benchmark Measures}
SVGEditBench evaluates changes to SVG code, whereas VectorGym covers generation from text and sketches, vector editing, and SVG captioning \citep{svgeditbench2024,vectorgym2026}.
MiLDEdit introduces editing of multi-layer documents with dedicated assessment axes \citep{mildedit2026}.
These studies refute an overly broad claim that there is no quantitative editing evaluation.
The remaining challenge is \emph{cross-representation comparability}: the permitted editing operations and document structures differ.

GraphicDesignBench probes multiple professional design capabilities in both understanding and generation settings \citep{deganutti2026gdb}.
BizGenEval evaluates commercial visual generation in five document domains with four capability axes and human-verified constraint checks \citep{li2026bizgeneval}.
AesEval-Bench focuses on whether visual-language models can assess aesthetic quality, including localized judgments \citep{an2026aeseval}.
These evaluations provide valuable but different evidence.
In particular, aesthetics-judge agreement, typography fidelity, valid code, and instruction-following should be reported separately.
The arXiv abstract and software repository for GraphicDesignBench currently disagree about task count (50 versus 49), so a definitive count requires version reconciliation.

\subsection{Protected Content Versus Necessary Reflow}
The advice to edit only the specified object is incomplete when the artifact encodes relationships.
InfoEdit studies infographics where an expanded text block or inserted element may force global layout \emph{reflow} \citep{yang2026infoedit}.
An evaluation can specify the intended target set $T$, permitted dependent changes $R(D,u,T)$, and explicitly protected content $P$:
\begin{equation}
\mathcal{U}=T\cup R(D,u,T),\qquad P\cap\mathcal{U}=\emptyset.
\label{eq:reflow}
\end{equation}
Here $\mathcal{U}$ denotes the allowed change scope.
A successful edit may require adjusting elements in $R$ while preserving the semantic content or geometry of $P$.
In a raster image, protected items may be regions or propositions rather than identifiable object IDs; in code-native documents they may be DOM nodes or layer dependencies.
This partition is a recommended evaluation abstraction, not a metric claimed by the source papers.

MPTc-Bench studies poster localization across markets, where typography and aesthetic conventions can change while a film's identity must remain stable \citep{lin2026mptc}.
It likewise demonstrates why minimal pixel deviation is an inadequate generic objective.

\subsection{Minimum Reporting Contract}
For comparable evaluations, studies should disclose input artifacts, allowed editing actions/tools, output and file validity requirements, the scope of target/dependent/protected content, the treatment of invalid outputs, and calibration of any human or VLM judge.
Successful rendering of a preview does not by itself show that the saved layered, vector, or code-native artifact remains valid.

\subsection{Domain-Specific Understanding and Scientific Validity}
DesignProbe examines visual-language-model understanding of graphic-design elements, style, and metaphor \citep{designprobe2024}.
Scientific illustration benchmarks introduce additional correctness constraints: SridBench assesses whether a generated graphic communicates scientific concepts faithfully \citep{sridbench2025}, and GENFIG1 tests graphical summary generation from research-paper content \citep{genfig12026}.
SciFlow-Bench reverse-parses generated scientific diagrams into graphs to measure structural recoverability rather than merely visual resemblance \citep{sciflowbench2026}.
These benchmarks evaluate different stages of a scientific communication pipeline, so high performance on recognition or aesthetics does not imply correct arrows, graphs, or native editable objects.

\subsection{Document-Scale Editing}
PPTArena measures in-place modifications of real presentation decks under natural-language requests, including structural and visual assessment \citep{pptarena2025}.
Unlike a screenshot comparison, an editing evaluation should inspect the original document tree, the requested change, and whether the revised file remains valid and manipulable.
This setting complements the multi-layer design-editing tasks of MiLDEdit and the infographic reflow scenarios of InfoEdit \citep{mildedit2026,yang2026infoedit}.
